\section{Material Interfaces} \label{section:material_interfaces}

\label{para:interface_section_intro} % ---------- %

Building on our foundational definitions, we now explore the key functional regions of a material interface, the
engineering properties that govern strength, adhesion, and stability, and the principal physical effects arising at
these junctions, electronic phenomena and energy-transport processes.

We then discuss design parameters for tailoring interfacial behaviour, the role of crystallographic orientation and
polymorph-driven reconstruction, and the generation and prediction of interfacial defects.

Finally, we highlight representative applications, from gate dielectric layers and transition metal dichalcogenides
(TMDC) heterojunctions to energy-storage interfaces, thermoelectric contacts, and optoelectronic devices, and conclude
with current active modelling challenges.

\label{para:interfacial_engineering_properties} % ---------- %

In materials engineering, interfaces govern a wide array of properties including adhesion, corrosion resistance,
mechanical strength, and diffusion kinetics\cite{Leitherer2023-ew, Butler2019-ws}.

The strength and ductility of composite systems or coatings often depend on interfacial bonding and lattice
compatibility\cite{Lorentzon2025-si, Raabe2020-ra}.

Tailoring interfacial chemistry or morphology, through methods such as surface functionalisation, alloying, or atomic
layer deposition, is often essential for achieving desired macroscopic performance metrics\cite{Chagarov2015-ol, Leitherer2023-ew, Butler2019-ws}.

\label{para:electronic_interface_effects} % ---------- %

In electronic devices, interfaces critically influence charge transport, band alignment, and contact resistance\cite{Schusteritsch2015-lw, Leung2020-tv, Low2025-on}.

Semiconductor heterojunctions, metal-semiconductor contacts, and dielectric boundaries are fundamental to the operation
of transistors, diodes, and capacitors\cite{Low2025-on, Khot2025-hh}.

Localised changes in atomic registry, roughness, or defect density at the interface can significantly alter tunnelling
behaviour, barrier heights, and mobility\cite{Gerber2023-yd, Khot2025-hh}.

Where interfacial dipoles, charge redistribution, and electrostatic reconstruction result in complex electronic
responses not captured by bulk models\cite{Alam2020-ci}.

\label{para:energy_interfaces} % ---------- %

Energy systems are similarly interface-dominated.

In lithium-ion batteries, for example, the solid-electrolyte interphase (SEI) acts as a metastable interfacial layer
that dictates both stability and charge transport\cite{Butler2019-ws}.

The performance and longevity of solid-state batteries, fuel cells, and thermoelectric devices can depend on the
engineered properties of interfaces, which affect ionic conductivity, catalytic activity, thermal transport, and
electronic insulation\cite{Leitherer2023-ew}.

In photovoltaics, interfacial band offsets between donor and acceptor layers, or between absorber materials and charge
transport layers, govern charge separation efficiency and open-circuit voltage\cite{Low2025-on}.

\label{para:interface_design_parameters} % ---------- %

Overall, interfaces are not incidental features but active design parameters.

They serve as sites of emergent phenomena, such as interfacial polarisation, bandgap renormalisation, and phonon
scattering, that cannot be predicted from bulk properties alone.

As such, mastering interfacial science is a gateway to developing high-performance materials and devices.

\label{para:interface_definition} % ---------- %

The structural features that influence interface behaviour include crystallographic orientation, the nature of the phase
boundary, and the presence of defects\cite{Zhang2016-lt}.

\label{para:orientation_effects} % ---------- %

The relative orientation of crystals on either side of an interface can significantly affect the atomic coordination,
bond continuity, and symmetry properties across the interface\cite{Kuhbach2021-gz, Zhang2016-lt}.

Epitaxial alignment, where lattice planes across the interface maintain a well-defined relationship, typically results
in low-strain, coherent interfaces\cite{}.

In contrast, a mismatch in orientation can lead to misfit dislocations or interfacial strain, influencing charge carrier
transport, interfacial states, and mechanical stability\cite{Gerber2023-yd, Leung2020-tv, Leitherer2023-ew}.

Orientation also affects the formation energy and electronic reconstruction at the interface, which can lead to
phenomena like interface-induced states or altered band alignments\cite{Alam2020-ci, Chagarov2015-ol, Zhang2016-lt}.

% ---------- %
